\documentclass[a4paper,twoside]{article}
\usepackage[T1]{fontenc}
\usepackage[bahasa]{babel}
\usepackage{graphicx}
\usepackage{graphics}
\usepackage{float}
\usepackage[cm]{fullpage}
\pagestyle{myheadings}
\usepackage{etoolbox}
\usepackage{setspace} 
\usepackage{lipsum} 
\setlength{\headsep}{30pt}
\usepackage[inner=2cm,outer=2.5cm,top=2.5cm,bottom=2cm]{geometry} %margin
% \pagestyle{empty}

\makeatletter
\renewcommand{\@maketitle} {\begin{center} {\LARGE \textbf{ \textsc{\@title}} \par} \bigskip {\large \textbf{\textsc{\@author}} }\end{center} }
\renewcommand{\thispagestyle}[1]{}
\markright{\textbf{\textsc{AIF184001/AIF184002 \textemdash Rencana Kerja Tugas Akhir \textemdash Sem. Ganjil 2026/2027}}}

\newcommand{\HRule}{\rule{\linewidth}{0.4mm}}
\renewcommand{\baselinestretch}{1}
\setlength{\parindent}{0 pt}
\setlength{\parskip}{6 pt}

\onehalfspacing

\begin{document}
	
	\title{\@judultopik}
	\author{\nama \textendash \@npm} 
	
	%tulis nama dan NPM anda di sini:
	\newcommand{\nama}{Bram Matthew Pasaribu}
	\newcommand{\@npm}{6182101043}
	\newcommand{\@judultopik}{Perangkat Lunak Pembuat Visualisasi Graf Situs Web} % Judul/topik anda
	\newcommand{\jumpemb}{1} % Jumlah pembimbing, 1 atau 2
	\newcommand{\tanggal}{18/02/2026}
	
	% Dokumen hasil template ini harus dicetak bolak-balik !!!!
	
	\maketitle
	
	\pagenumbering{arabic}
	
	\section{Deskripsi}
	Tugas akhir ini bertujuan untuk merancang dan mengembangkan sebuah perangkat lunak yang mampu melakukan visualisasi struktur graf dari sebuah situs web. Perangkat lunak ini bekerja dengan cara mengambil dokumen \textit{Hypertext Markup Language} (HTML) dari suatu alamat situs web, untuk dilakukan proses \textit{parsing} sehingga membangun struktur \textit{Document Object Model} (DOM). Struktur DOM tersebut selanjutnya direpresentasikan dalam bentuk graf menggunakan \textit{Graphviz} yang menunjukkan hubungan hierarki antar elemen HTML.

    Alur kerja sistem dimulai dari pengguna yang memasukkan \textit{Uniform Resource Locator} (URL). Sistem kemudian mengambil dokumen HTML melalui proses \textit{web request} dan melakukan \textit{parsing} untuk membentuk struktur DOM. Selanjutnya, sistem menelusuri struktur DOM menggunakan metode \textit{Depth-First Search} (DFS), di mana setiap elemen direpresentasikan sebagai simpul (\textit{node}) dan hubungan induk-anak (\textit{parent-child}) sebagai sisi (\textit{edge}).
    
    Hasil dari proses tersebut disimpan dalam format .dot yang kemudian divisualisasikan menggunakan \textit{Graphviz} menjadi bentuk graf. Visualisasi ini membantu pengguna dalam memahami struktur halaman web secara lebih jelas.
	
	\section{Rumusan Masalah}
	\begin{enumerate}
        \item Bagaimana cara mengambil dokumen \textit{HTML} dari suatu situs web berdasarkan \textit{URL} serta membangun struktur \textit{Document Object Model} (DOM)?
        
        \item Bagaimana cara merepresentasikan struktur \textit{DOM} ke dalam bentuk graf yang terdiri dari simpul (\textit{node}) dan sisi (\textit{edge}) melalui proses traversal yang sistematis?
        
        \item Bagaimana merancang perangkat lunak yang mampu menghasilkan visualisasi graf dari struktur halaman web secara otomatis dan mudah dipahami pengguna?
    \end{enumerate}
	
	\section{Tujuan}
	\begin{enumerate}
    \item Mengambil dokumen \textit{HTML} dari suatu situs web berdasarkan \textit{URL} serta membangun struktur \textit{Document Object Model} (DOM).
    
    \item Merepresentasikan struktur \textit{DOM} ke dalam bentuk graf yang terdiri dari simpul (\textit{node}) dan sisi (\textit{edge}) melalui proses traversal yang sistematis.
    
    \item Merancang perangkat lunak yang mampu menghasilkan visualisasi graf dari struktur halaman web secara otomatis dan mudah dipahami pengguna.
    \end{enumerate}
	
	\section{Deskripsi Perangkat Lunak}
	Perangkat lunak yang akan dibuat memiliki fitur utama sebagai berikut:
    \begin{itemize}
    \item Menerima masukan berupa \textit{Uniform Resource Locator} (URL) dan mengambil dokumen \textit{Hypertext Markup Language} (HTML) dari situs web yang diberikan pengguna.
    
    \item Melakukan proses \textit{parsing} untuk membangun struktur \textit{Document Object Model} (DOM) serta menelusuri struktur tersebut untuk mengidentifikasi hubungan antar elemen.
    
    \item Merepresentasikan struktur \textit{DOM} ke dalam bentuk graf berupa simpul (\textit{node}) dan sisi (\textit{edge}), serta menghasilkan visualisasi graf yang dapat membantu pengguna memahami struktur halaman web.
    \end{itemize}
	
	\section{Detail Pengerjaan Tugas Akhir}
	Bagian-bagian pengerjaan tugas akhir ini adalah sebagai berikut:
    \begin{enumerate}
        \item Melakukan studi literatur mengenai HTML, DOM, graf, dan metode traversal.
        \item Mempelajari teknik \textit{parsing} HTML serta visualisasi graf.
        \item Menganalisis kebutuhan dan merancang sistem perangkat lunak.
        \item Mengimplementasikan sistem untuk mengambil HTML, membangun DOM, dan membentuk graf.
        \item Melakukan pengujian dan evaluasi sistem.
        \item Menyusun dokumentasi dan laporan tugas akhir.
    \end{enumerate}
	
	\section{Rencana Kerja}
	Rincian capaian yang direncanakan sebelum Tugas Akhir 1 adalah sebagai berikut:
    \begin{enumerate}
        \item Melakukan studi literatur mengenai HTML, DOM, graf, dan metode traversal.
        \item Mempelajari teknik \textit{parsing} HTML dan visualisasi graf.
        \item Menganalisis kebutuhan serta merancang sistem perangkat lunak.
        \item Menulis sebagian dokumen tugas akhir (Bab 1, 2, dan 3).
    \end{enumerate}
    
    Sedangkan yang akan diselesaikan di Tugas Akhir 1 dan Tugas Akhir 2 adalah sebagai berikut:
    \begin{enumerate}
        \item Mengimplementasikan sistem pengambilan HTML, parsing DOM, dan pembentukan graf.
        \item Melakukan pengujian dan evaluasi hasil visualisasi.
        \item Menyelesaikan penulisan dokumen tugas akhir (Bab 4, 5, dan 6).
    \end{enumerate}
	
	\vspace{1cm}
	\centering Bandung, \tanggal\\
	\vspace{2cm} \nama \\ 
	\vspace{1cm}
	
	Menyetujui, \\
	\ifdefstring{\jumpemb}{2}{
		\vspace{1.5cm}
		\begin{centering} Menyetujui,\\ \end{centering} \vspace{0.75cm}
		\begin{minipage}[b]{0.45\linewidth}
			% \centering Bandung, \makebox[0.5cm]{\hrulefill}/\makebox[0.5cm]{\hrulefill}/2013 \\
			\vspace{2cm} Nama: \makebox[3cm]{\hrulefill}\\ Pembimbing Utama
		\end{minipage} \hspace{0.5cm}
		\begin{minipage}[b]{0.45\linewidth}
			% \centering Bandung, \makebox[0.5cm]{\hrulefill}/\makebox[0.5cm]{\hrulefill}/2013\\
			\vspace{2cm} Nama: \makebox[3cm]{\hrulefill}\\ Pembimbing Pendamping
		\end{minipage}
		\vspace{0.5cm}
	}{
		% \centering Bandung, \makebox[0.5cm]{\hrulefill}/\makebox[0.5cm]{\hrulefill}/2013\\
		\vspace{2cm} Nama: \makebox[3cm]{\hrulefill}\\ Pembimbing Tunggal
	}
\end{document}
